\section{A modular transfer of actual root-of-unity poles}
\label{sec:modular-pole-transfer}

The transfer in this section concerns ordinary Hilbert denominators.
It uses a congruence of the actual coefficient sequences and a
specific noncancellation statement after reduction modulo three.
No equality between a universal denominator and a Hilbert
denominator is assumed.

\begin{lemma}[Cyclotomic factors after reduction]
Let $p$ be a prime and write $d=p^r e$, where $p\nmid e$.
In $\mathbb F_p[t]$ one has
\[
 \overline{\Phi_d}(t)=
 \begin{cases}
 \overline{\Phi_e}(t),&r=0,\\
 \overline{\Phi_e}(t)^{p^{r-1}(p-1)},&r\ge1.
 \end{cases}
 \tag{PNEW-M1}
\]
The roots of $\overline{\Phi_e}$ in
$\overline{\mathbb F}_p$ are exactly the elements of
multiplicative order $e$, and each is simple.
\end{lemma}
\begin{proof}
For $p\nmid e$, the characteristic-zero cyclotomic identity is
\[
 \Phi_{ep^r}(t)\Phi_e(t^{p^{r-1}})=\Phi_e(t^{p^r})
 \qquad(r\ge1).
\]
For example, it follows by grouping the roots on the right
according to whether their order has $p$-part $p^r$ or at
most $p^{r-1}$. Reduction modulo $p$ replaces each substitution
$t\mapsto t^{p^j}$ by the $p^j$th power. Cancelling the nonzero
polynomial $\overline{\Phi_e}^{\,p^{r-1}}$ in the integral
domain $\mathbb F_p[t]$ proves (PNEW-M1).

The factorization $t^e-1=\prod_{b\mid e}\Phi_b(t)$ remains
valid modulo $p$, and $t^e-1$ is separable because its derivative
is $e t^{e-1}$. Induction on $e$ in that factorization assigns
to $\overline{\Phi_e}$ precisely those roots whose exact
order is $e$. Separability proves the final assertion.
\end{proof}

\begin{lemma}[An isolated graded fixed point gives no pole]
Let $R=\bigoplus_{j\ge0}R_j$ be a finitely generated graded
$\mathbb C$-algebra with $R_0=\mathbb C$ and finite-dimensional
pieces. If the fixed locus of its grading subgroup $\mu_d$ is
supported only at the graded origin, then $H_R(t)$ is regular
at every primitive $d$th root of unity.
\end{lemma}
\begin{proof}
The fixed scheme is defined by the homogeneous ideal
\[
 I_d=(R_j:d\nmid j).
\]
Indeed a homogeneous function of degree $j$ has pullback
eigenvalue $\zeta_d^j$, so these are exactly the equations
of the grading fixed scheme. Choose finitely many homogeneous
generators $f_1,\ldots,f_b$ of this ideal from those pieces,
and put $a_i=\deg f_i$, so $d\nmid a_i$.
The hypothesis says that $R/I_d$ is finite-dimensional.

Every homology module of the Koszul complex on the $f_i$ is
finitely generated over $R$ and annihilated by each $f_i$:
multiplication by a Koszul generator is chain homotopic to zero.
It is therefore finite-dimensional. Taking the graded Euler
characteristic of that complex yields
\[
 H_R(t)\prod_{i=1}^b(1-t^{a_i})
       =\sum_{i=0}^b(-1)^i H_{H_i(f_1,\ldots,f_b;R)}(t),
 \tag{PNEW-M2}
\]
whose right side is a polynomial. None of the factors on the
left vanishes at a primitive $d$th root. This proves regularity.
Nilpotents supported at the origin do not affect the argument.
\end{proof}

\begin{theorem}[PNEW-013: modular pole transfer]
Let $F,G\in\mathbb Q(t)\cap\mathbb Z[[t]]$ have all their
finite poles at roots of unity. Let $p$ be a prime, let
$p\nmid M$, and suppose
\[
                  F(t)\equiv G(t)\pmod{p\mathbb Z[[t]]}.
 \tag{PNEW-M3}
\]
Suppose $G$ has a pole of order $\delta\ge1$ at a primitive
$M$th root $\zeta$, has no pole at any primitive
$Mp^r$th root for $r\ge1$, and
\[
 C=\lim_{t\longrightarrow\zeta}(1-t^M)^\delta G(t)
 \tag{PNEW-M4}
\]
is a unit in $\mathbb Z[\zeta]_{\mathfrak p}$ for at least
one prime $\mathfrak p$ over $p$.
Then $F$ has an actual pole at a primitive $Mp^r$th root for
some $r\ge0$. In particular every presentation
\[
                  F(t)=\frac{P(t)}{\prod_i(1-t^{a_i})^{b_i}}
 \qquad(a_i,b_i\in\mathbb Z_{>0},\ P\in\mathbb Q[t])
 \tag{PNEW-M5}
\]
has $a_i\ge M$ for at least one $i$.
\end{theorem}
\begin{proof}
A reduced rational function of the stated kind admits a
denominator which, up to sign, is a finite product of
cyclotomic polynomials. Normalize its constant coefficient
to one. Its numerator is integral: multiplication of its
integral formal series by that integral denominator gives
an integral formal series which is a polynomial. Thus write
\[
 F=N_F/D_F,\qquad G=N_G/D_G,
 \qquad N_F,N_G,D_F,D_G\in\mathbb Z[t],
 \quad D_F(0)=D_G(0)=1,
 \tag{PNEW-M6}
\]
with the pairs coprime over $\mathbb Q[t]$ and both denominators
products of signed cyclotomic factors.

Write
$D_G=\varepsilon\Phi_M^\delta D_0$, with $\varepsilon\in\{1,-1\}$.
The no-higher-pole hypothesis and (PNEW-M1) imply that
$\overline{D_0}$ is nonzero at
$\alpha=\zeta\bmod\mathfrak p$. Here $\alpha$ has exact
order $M$, again by (PNEW-M1), since $p\nmid M$.
Also put $U_M=(1-t^M)/\Phi_M(t)$.
Separability of $t^M-1$ modulo $p$ gives
$\overline{U_M}(\alpha)\ne0$.
Consequently $D_0(\zeta)$ and $U_M(\zeta)$ are units at
$\mathfrak p$, and (PNEW-M4) becomes
\[
                  C=
 \frac{U_M(\zeta)^\delta N_G(\zeta)}
      {\varepsilon D_0(\zeta)}.
 \tag{PNEW-M7}
\]
The assumed unit property forces $N_G(\zeta)$ to be a unit.
Thus $\overline{N_G}(\alpha)\ne0$, whereas
$\overline{D_G}$ vanishes there to order exactly $\delta$.
The rational function $\overline{G}\in\mathbb F_p(t)$ has
a genuine pole at $\alpha$.

Reduction of either expression in (PNEW-M6) is well defined
both in $\mathbb F_p(t)$ and in $\mathbb F_p[[t]]$, because
each denominator has constant coefficient one. The natural
map from such rational functions to formal series is
injective. Equation (PNEW-M3) therefore identifies
$\overline F$ and $\overline G$ as rational functions.
It follows that $\overline{D_F}(\alpha)=0$.

Some actual denominator factor $\Phi_d$ of $D_F$ consequently
has $\overline{\Phi_d}(\alpha)=0$. By (PNEW-M1), the part
of $d$ prime to $p$ must equal the exact order $M$ of
$\alpha$. Hence $d=Mp^r$ for some $r\ge0$.
Because $N_F,D_F$ were coprime in characteristic zero,
this is an actual pole of $F$. Rationality gives the same
pole order at every primitive $d$th root.

Finally, the denominator in (PNEW-M5) must vanish at any
such root. Some $a_i$ is divisible by $d=Mp^r$, and
therefore $a_i\ge M$. This last implication does not require
positivity of $P$.
\end{proof}

\paragraph{Why the higher-lift hypothesis is necessary.}
One cannot replace the hypotheses by a unit residue alone.
For instance
\[
 G(t)=1+\frac{p\,t^{pM}}{1-t^{pM}}
 \tag{PNEW-M8}
\]
has integral nonnegative coefficients, constant term one,
and $\lim_{t\to\zeta}(1-t^M)G(t)=1$.
But $\overline G=1$ has no pole. The additional primitive
$pM$th poles coalesce with the primitive $M$th roots after
reduction, and their denominator derivative supplies the
factor $p$ which invalidates the naive residue inference.

Nor does the theorem force the exact order $M$ for $F$.
For $p=3$, $M=2^a\ge4$ and $z=t^{M/2}$, put
\[
 F=\frac{1+z}{1-z+z^2}+\frac{2t}{(1-t)^2},\qquad
 G=\frac1{1+z}+\frac{2t}{(1-t)^2}.
 \tag{PNEW-M10}
\]
Their difference is $3z/(1+z^3)$, so they have congruent
integral coefficients. Both have constant term one and
nonnegative coefficients: the first summands of $F$ have
successive $z$-coefficients $(1,2,1,-1,-2,-1)$ periodically,
and those of $G$ alternate between one and minus one;
the added term contributes $2j$ in every positive $t$-degree
$j$. The primitive-$M$ residue of $G$ is two, a three-adic
unit, but $F$ has its corresponding nontrivial poles only
at order $3M$, since $1-z+z^2=\Phi_{3M}(t)$.
This demonstrates the exact-order limitation even for proper
rational functions with nonnegative integral coefficients.
These examples are not asserted to be invariant-ring Hilbert
series.

\begin{corollary}[Unconditional tensor consequence: numerical PH3 fails]
For every $n\ge2$, there exists an integer $r_n\ge0$ such that
$H_n^{\mathrm{ten}}$ has an actual pole at every primitive
$2^{n-1}3^{r_n}$th root of unity. Every rational presentation
of this series with denominator a product of factors
$(1-t^{a_i})^{b_i}$ therefore has $\max_i a_i\ge2^{n-1}$.
In particular the existence clause of GCT VI,
Hypothesis1.6.7, fails for triple-rectangular Kronecker
coefficients. In this realization the unified parameter of
GCT VI, Section3.4, is $\operatorname{rank}(x)=n^2+5n$;
the introductory parameter
$h=\dim G'+\operatorname{ht}(\lambda)+\operatorname{ht}(\pi)$
is instead $n^4+O(n)$. The exponential lower bound exceeds
every polynomial in either parameter.
\end{corollary}
\begin{proof}
Put $M=2^{n-1}$. K-023 identifies the cyclic-permutation
trace series $J_n$ with the weighted $SL_n$-invariant Hilbert
series of $W\oplus\mathbb S_{(2,1)}W$, with summand weights
one and three and total weight divided by $n$.
K-024, using the all-dimensional Lie proof R-NEW-004 and
the fully checked support theorem P-F30, proves that its
nonzero $\mu_M$-fixed quotient locus is one grading orbit,
with full grading stabilizer exactly $\mu_M$, density-one
curve closure, and smooth ambient quotient along that orbit.
Its nontrivial moving eigenvalues are powers of a primitive
$M$th root. These are proved inputs, not additional
hypotheses of this corollary.

For clarity, the coefficient congruence can be checked without
using the cycle-index identity itself. On the rational
multiplicity space
\[
 (([s^n]_{\mathbb Q})^{\otimes3})^{S_{ns}},
\]
the cyclic factor permutation has order three. Every rational
representation of this cyclic group is a sum of trivial lines
and copies of the two-dimensional rational representation
with eigenvalues $\omega,\omega^2$.
If their multiplicities are $a,b$, its dimension is $a+2b$
and its trace is $a-b$. Thus
\[
                     k_n(s)\equiv[t^s]J_n(t)\pmod3
                         \qquad(s\ge0).
 \tag{PNEW-M9}
\]
The Hilbert series are rational with root-of-unity poles,
because the invariant rings are finitely generated and
positively graded. Their coefficients are integers.

By these geometric results and the proved local
Hilbert formula P-F23, $J_n$ has a simple primitive-$M$ pole
whose coefficient is the inverse product of its moving
factors $1-\zeta^j$, with $M\nmid j$.
Every such factor is a unit at every prime above three:
the reduction of $\zeta$ retains exact order $M$.
The density is one, so its leading coefficient is a
three-adic unit.

For $r\ge1$, the $\mu_{M3^r}$ fixed locus is contained
in the $\mu_M$ fixed locus. No nonzero point of the latter
orbit is fixed, because its grading stabilizer is exactly
$\mu_M$. Thus the larger fixed locus is supported only
at the origin. The preceding Koszul lemma excludes all
primitive-$M3^r$ poles of $J_n$.
Theorem PNEW-013 now applies with $p=3$, $F=H_n^{\mathrm{ten}}$
and $G=J_n$.
\end{proof}

\paragraph{An explicit finite range for the higher lift.}
Combining the preceding corollary with the already proved uniform
lift bound K-009 gives
\[
 2^{n-1}3^{r_n}\le
       \frac{(3-2/n)^{(3n-2)/2}}{n}.
 \tag{PNEW-M11}
\]
Equivalently, $0\le r_n\le R_n$, where $R_n$ is the largest
nonnegative integer $r$ satisfying the entirely integral test
\[
       4^{n-1}3^{2r}n^{3n}\le(3n-2)^{3n-2}.
 \tag{PNEW-M12}
\]
Indeed K-009 excludes every actual pole whose order exceeds
its displayed bound, so it applies directly to the pole supplied
by the corollary. For example, $R_2=R_3=0$,
$R_4=R_5=1$, $R_9=4$, $R_{26}=18$, and $R_{58}=45$.
In particular the possible exponent grows at most linearly with
$n$. This is a corollary of K-009 and PNEW-013, not a new
determinant bound. It neither chooses the actual exponent nor
gives an archimedean lower bound for its wave amplitude.

The modular index of every positive rational form for
$H_n^{\mathrm{ten}}$ is therefore at least $2^{n-1}$ for
every $n\ge2$, so no polynomial in either
$\operatorname{rank}(x)=n^2+5n$ or the introductory
$h=n^4+O(n)$ bounds it.
This is an unconditional refutation of the existence clause
of numerical PH3. It does not identify the
integer $r$, prove a pole at the exact primitive order
$2^{n-1}$, or bound the strict positivity and saturation
indices. Those are distinct questions. For $n=3,4,5$, the
separate exact calculations K-001, K-008 and K-014 establish
the pole with $r=0$; no such uniform exact-order conclusion
is inferred here.
